\section{The sharp universal Euler constant}
\label{ue:sec:universal}
The Gaussian axis maximum becomes the sharp Euler constant only after a
separate theorem controls all later scaled errors. The following result
supplies that missing implication; complete proofs and exact rational
certificates accompany \cite{UniversalEulerReport}.

\begin{theorem}[Universal real-order Euler bound]\label{ue:thm:universal}
For $a,b>0$, let $g_{a,b}=\Imag\Li_{a,b}(\mathrm i,1)$ have its principal
analytic meaning, and set
\[
 A_k=H_{2k}^{(b)}/(2k+1)^a,\qquad
 E_N=\sum_{j=0}^{N-1}2^{-j-1}\sum_{k=0}^j(-1)^k\binom jk A_k.
\]
Then
\[
 0<E_N-g_{a,b}<\frac98\,2^{-N}\qquad(N\ge2).
\]
The least constant uniform in $a,b>0$ and $N\ge1$ is
\[
 C_* = \max_{b>0}\{\beta(b)+2^{-b}\eta(b)\},\qquad
 1.1365611033<C_*<1.1365611046.
\]
Its supremum is approached at $N=1$, $a\downarrow0$, and the unique
maximizing inner order $b_*$, which satisfies $1.3021<b_*<1.3023$.
\end{theorem}
\begin{proof}
For independent rate-one gamma variables $S,T$ of shapes $a,b$, put
$X=e^{-S}$, $Y=e^{-T}$, $p=X^2$. Gamma moments and finite geometric summation give
\[
 2^N(E_N-g_{a,b})=\mathbb E\left[
 \frac{f_N(pY^2)-f_N(p)}{1-Y}\right],\qquad
 f_N(v)=\frac{(1-v)^N}{1+v}.
\]
This is strictly positive, without any assumption on convergence of the
original boundary series. The rational kernel is strictly less than
$9/8$ for every $N\ge2$ and $0\le p,Y\le1$ with continuous endpoint
interpretation. For $N=2,3$ this follows from two polynomial Bernstein
certificates. For $N\ge4$, expand $1/(1+v)$ geometrically and bound the
binomial divided differences by
\[
 (1+Y)\left(\frac{1-Y^2}{1-Y^{2m/(m-1)}}\right)^{m-1}.
\]
These envelopes decrease in $m$; the fourth is below $9/8$ by a third
polynomial certificate. The report supplies all 1542 positive rational
coefficients and two independent exact replay routes.

For $N=1$ the kernel is
$2p(1+Y)/((1+p)(1+pY^2))$, strictly increasing in $p$. Its expectation
is strictly below the axis value $\beta(b)+2^{-b}\eta(b)$ for $a>0$
and tends to it as $a\downarrow0$. The inherited unique Gaussian maximum
therefore controls $N=1$. Since $\pi/4+\log(2)/2>9/8$, every later
truncation is strictly below that maximum. Exact finite Euler intervals
and an analytic curvature bound give the stated rational enclosures.
\end{proof}

The restricted constant-one and sharp critical/subcritical bounds remain
stronger on their respective domains. Scaled errors are not generally
monotone: at $(a,b)=(4,1)$ one has
$2^2(E_2-g)>2(E_1-g)$. Thus the later-truncation theorem cannot be replaced
by an unproved monotonicity assertion. The global integer normalized-radius
and mixed $S_6$ conjectures are unaffected.
